% Problema 4, parte 5: campi da incollare nella piattaforma (uno per sezione)

% ===== CAMPO: 1. Result and scope =====
\textbf{Official request.} \#\# Parte 5 (C5) --- The certified boundary

\textbf{Punteggio:} 8 points · \textbf{Valutazione:} Judged

The main certification cell: an initial range of sizes as long as you
can make it, plus two isolated sizes further out. Determine the largest
\(k\) for which you can certify the statement for all sizes up to and
including \(k\), and certify it. The range you claim must be contiguous,
and if your argument at a given size assumes that all smaller sizes have
already been settled you must say so. Then decide the two isolated sizes
\(k = 24\) and \(k = 30\): show that neither is the least size at which
the statement can fail.

For this cell the exhaustiveness certificate of the hand-in rules is not
enough on its own. Add:

\begin{enumerate}
\def\labelenumi{(\roman{enumi})}
\item
  a demonstration that your method is not vacuous. Construct a case in
  which an admissible configuration genuinely exists, run your machinery
  on it unmodified, and show it returns that configuration. A method
  that reports ``nothing survives'' at every size it is pointed at is
  indistinguishable from a method with a bug, and will be graded as one;
\item
  for every object your pruning leaves undecided, at every \(k\) you
  claim --- not a sample --- an individual decision, plus the smallest
  part of that object which already forces the decision, plus a proof
  that no smaller part does;
\item
  the exact list, not merely the count, of what survives at each \(k\);
\item
  every pruning rule you use beyond those you have proved, proved. If
  your search needs a rule you invented to finish a size, that rule is
  part of your claim: state it, prove it, and show the survivor list is
  unchanged when you switch it off. A size that only closes with an
  unproved rule is not certified;
\item
  a second implementation, written independently of your first, that
  differs in method --- not the same algorithm twice --- and the two
  survivor lists compared elementwise at every \(k\) you claim. Report
  any difference rather than reconciling it silently: a disagreement
  means one of them is wrong, and finding which is part of the cell.
\end{enumerate}

\textbf{What we deliver.} No certificate. Observation: for \(k=24,30\)
brute force on \(L_k\) is impractical; a size-reduction argument is
needed (if a counterexample exists at \(k\), one exists at some
\(k'<k\)).


% ===== CAMPO: 2. Proof =====
Not developed.


% ===== CAMPO: 3. Verification: instructions, dependencies, timings =====
Available code (Python 3, standard library; each script runs in under a
minute): - No code yet.


% ===== CAMPO: 4. Sources and contribution =====
arXiv literature (deterministic search
\texttt{tools/cerca\_letteratura.sh}, abstracts read, not used as
proof): - arXiv:2603.26043v1 --- Finiteness of Disjoint Covering Systems
with Precisely One Repeated Modulus (Yu Hashimoto, 2026); abstract only
read. - arXiv:1511.04293v1 --- Searching for Disjoint Covering Systems
with Precisely One Repeated Modulus (Shalosh B. Ekhad, Aviezri S.
Fraenkel, Doron Zeilberger, 2015); abstract only read. -
arXiv:2607.24655v1 --- On the problem of large gcd for disjoint residue
classes (Jan Fornal, Yu-Chen Sun, 2026); abstract only read. -
arXiv:2608.15873v1 --- Two Questions on \(G\)-harmonic Tuples (Murali
Menon, 2026); abstract only read. As in part 3.


% ===== CAMPO: 5. Limits and unresolved parts =====
Everything.

